\section{Two-sided extensions of a coefficient-jet core}

\subsection*{Finding M-019: a two-ideal support criterion for full multiplicities}
Let \(r\geq2\), \(\ell\geq0\), and define
\[
 R=r-1,\quad b=2\ell+2,\quad h=\lceil\sqrt r\rceil,
 \quad Q=\binom{R+\ell}{\ell},\quad T=hQ.
\]
Let \(A,B\subset\mathbb Z_{\geq0}^{R}\) be nonempty
finite monomial order ideals. Suppose that a finite coefficient
support satisfies
\[
                         \Gamma\subseteq A+B+D_b,
      \qquad D_b=\{\eta\in\mathbb Z_{\geq0}^{R}:|\eta|\leq b\}.
                                                               \tag{M-019.1}
\]
Then the determinant coefficient projection onto \(\Gamma\)
is dominant whenever
\[
                            n\geq N=T(|A|+|B|).          \tag{M-019.2}
\]
The internally padded permanent has the same dominant projection
whenever \(n\geq m\geq N\). Both differentials have explicit
elementary unit minors over \(\mathbb Z\).

Consequently, if \(\lambda\vdash nd\), \(\ell(\lambda)\leq r\),
and its full nested support \(\Gamma_\lambda\) from M-018
satisfies (M-019.1), then (M-019.2) gives
\(K_n(\lambda)=a_\lambda(d[n])\). Under the permanent-size
hypothesis, \(q^m_\lambda=K_n=a_\lambda\).
These assertions retain every highest-weight copy.

\paragraph{The constant bridge.}
Use \(h\) inward and \(h\) outward copies of the monomial
tree on \(R\) variables \(z_1,\ldots,z_R\), through degree
\(\ell\), just as in M-013. Each copy has \(Q\) vertices.
The inward copies point toward their roots, and the outward
copies point away from their roots. Since \(h^2\geq r=R+1\),
inject the \(R+1\) labels
\[
                             x,z_1,\ldots,z_R
\]
into distinct ordered pairs of an inward root and an outward
root, and add the corresponding directed bridges. Unused
pairs have no edge. The label \(x\) becomes one after
dehomogenization; it is still a linear homogeneous matrix
entry before that step. The graph remains acyclic.

Every monomial of degree between zero and \(b-1=2\ell+1\)
is now a pure inward-to-outward path product after \(x=1\).
For degree zero use the \(x\)-bridge with its two roots.
For positive degree, choose one factor for a variable bridge
and split the remaining factors into two monomials of degrees
at most \(\ell\). Their inward and outward tree paths,
with that bridge, give exactly the required monomial.
For fixed endpoints the path is unique: it must use the
single bridge between their two components. In particular
the added constant bridge introduces no alternate path.

\paragraph{The two attached order ideals.}
At each of the \(T\) old inward vertices, attach a separate
tree indexed by \(A\), identifying its zero vertex with
that old vertex and pointing every new edge toward it.
At each of the \(T\) old outward vertices, attach a separate
tree indexed by \(B\), again identifying the zero vertex,
but pointing every new edge away from the root.
For every nonzero multi-index, choose a supported variable
to remove for its parent; the order-ideal property ensures
that parent exists. Labels are the removed variables.
All nonroot attached vertices are disjoint.

The total size is
\[
        2T+T(|A|-1)+T(|B|-1)=T(|A|+|B|)=N.
\]
A path from an attached inward vertex \(\alpha\) to an
attached outward vertex \(\beta\), using an old selected
path of weight \(z^\eta\), is unique and has weight
\[
                              z^{\alpha+\eta+\beta}
                                                               \tag{M-019.3}
\]
after \(x=1\). No old edge enters a new inward nonroot
vertex, and no new outward nonroot vertex leads back into
the old graph. Thus the path must have precisely these
three segments. There are no directed cycles.

\begin{proof}[Dominance, fixed padding and all copies]
Let \(\mathcal B\) be the homogeneous adjacency matrix.
At the base pencils \(xI-\mathcal B\) and
\(xI+\mathcal B\), acyclicity gives determinant and
permanent both equal to \(x^N\). In the determinant,
the reverse-entry cofactor of a directed path is its
positive monomial product: its edge signs cancel the
cycle permutation sign. In the permanent the open-path
minor has the same product with positive sign. There
are no disjoint nontrivial cycles. These statements
also apply to a path using the constant \(x\)-bridge.

For \(\gamma\in\Gamma\), choose
\(\gamma=\alpha+\beta+\eta\) as in (M-019.1).
If \(|\eta|\leq b-1\), use its selected core path
and multiply the cofactor (M-019.3) by the linear entry
variation \(H=x\). If \(|\eta|=b\), remove one variable
for \(H\) and use the core path of the remaining degree
\(b-1\). Each resulting derivative is exactly
\(z^\gamma\) after \(x=1\), with no extra monomials.
This includes \(\eta=0\), which is precisely the purpose
of the constant bridge. Different target monomials use
different elementary source coordinates, so the selected
Jacobian minor is the identity over \(\mathbb Z\).

The degree bound causes no extra hypothesis. Attach a
maximum-degree member of \(A\) at the source of an old
degree-\(b-1\) path, and a maximum-degree member of \(B\)
at its target. The resulting simple path has
\(\max\deg A+(b-1)+\max\deg B\) edges. Therefore
\[
                       N\geq\max\deg A+\max\deg B+b,
\]
so all target coefficients are valid at degree \(N\).
The degree-zero core path has one edge rather than zero,
which only makes this bound easier in that endpoint case.
For larger matrices append \(xI\) blocks.

Restricting to the selected source coefficient directions
and translating the target by its base value gives
\(F(t)=t+O(t^2)\) over \(\mathbb Z\). The least
homogeneous term of a nonzero polynomial survives pullback.
This proves injection of the polynomial coordinate algebra
after every commutative base change, and in particular
dominance over \(\mathbb C\).

Every selected derivative varies an entry with one inward
and one outward endpoint. These vertex sets are disjoint,
even for the degree-zero core path. Thus all selected
directions are off diagonal. Hold any one matrix diagonal
fixed at \(x\), and simultaneously permute rows and
columns to make it the distinguished entry \(X_{11}\).
The resulting internally padded permanent is exactly
\(x^{n-m}\operatorname{per}(xI+\mathcal B)\) with the
allowed off-diagonal variations. Its dehomogenized detected
coefficients are unchanged. This proves the permanent
claim with the stated internal-padding convention.

All substituted forms belong to the respective orbit
closures by endomorphism substitution. F-016 places every
highest-weight ideal vector of type \(\lambda\) in the
algebra generated by \(\Gamma_\lambda\). The proved
coordinate injection forces such an ideal vector to be
zero. Hence the corresponding entire ideal multiplicity
space vanishes, giving the asserted ambient equalities.
\end{proof}

\paragraph{A useful two-prefix consequence in BIP range.}
For every actual obstruction at \(n\geq m^{25}\),
\[
                   \lambda_4\geq27
                       \quad\text{or}\quad\lambda_2\geq239063403.
                                                               \tag{M-019.4}
\]
Indeed use \(b=26\), and suppose
\(\lambda_2\leq26+t\), \(\lambda_4\leq26\).
Every target monomial can have all but at most twenty-six
factors removed from the first two nonfirst variables.
Its removed part has exponents at most \(t\) in each.
Thus (M-019.1) holds with
\[
             A=\{1,x_2,\ldots,x_2^t\},\qquad
             B=\{1,x_3,\ldots,x_3^t\}.
\]
The known permanent-size reduction gives \(m\geq190\).
Since \(r\leq m^2\), the new value
\(h=\lceil\sqrt r\rceil\) is still at most \(m\).
Hence \(N\leq2(t+1)m\binom{m^2+11}{12}\).
M-017's exact constant \(C_*=478126755\) shows this
fits whenever \(2(t+1)\leq C_*\).
The largest such integer is \(t=239063376\), proving
(M-019.4). The low-length cases were already excluded by
O-006--O-007; the argument uses \(r\geq4\) otherwise.
The general theorem also permits capped or nested ideals
in each direction, rather than only these univariate ones.

\paragraph{Geometric dictionary and scope.}
The two order ideals and the jet core index finite free
torus-weight coefficient modules. Their bridge and tree
matrices are integral, and the unit differential minor
gives actual coordinate-algebra injection under every
base change. The torus grading and ideal comodules are
SGA~3, I, Propositions~4.7.2--4.7.3, pp.~34--36.
F-016 realizes the relevant Schur operations through
SGA~6, VI, Theorem~3.3, pp.~392--393, and Macdonald,
I(4.3), p.~63, Appendix~A(5.3)--(5.6), pp.~155--156.
Projection dominance compares all multiplicity copies;
it does not assert exact representation of all forms with
the target support. The explicit strictness example O-025
separates this two-sided criterion from every one-sided
order-ideal construction of M-018. No general saturation
or actual multiplicity obstruction is asserted.
